\scanpage{161}
\mspage{80}
% Lignes 1--5 (continuation de la note marginale de la page précédente).
\illegible{2}, et \unsure{au} \deleted{\unsure{le}} fait qu'elles restent \inserted{très} proches du langage \unsure{et|\deleted{et}} de l'intuition topologique, elle est supérieure au $p^t$ de vue ``groupes à lacets'', qui correspond à l'appr. du calculatoire plutôt.)

% Lignes 5--16 : grand passage encadré et barré par des traits obliques.
\deleted{, donnons-nous \deleted{une partie sous topos} (\inserted{\unsure{néc$^t$.}}\unsure{ouvert}) $B_{D_1^*}$ de $B_{D^*}$ (correspondant à une partie $I_1$ de $I=\pi_0(B_{D^*})$) de sorte que
\[
 B_{D^*}\simeq B_{D_0^*}\amalg B_{D_1^*}
\]
\margin{($D_0^*$ correspond à la partie $I_0=I\setminus I_1$ de $I$)}
On a une décomposition analogue
\[
 B_D\simeq B_{D_0}\amalg B_{D_1}.
\]
\inserted{\unsure{déf.} $B_{\pi_0(D^*)}=B_S$}
Considérons le topos \unsure{$B_U'$} ``somme amalgamée'' de \deleted{\unsure{$B_U$}}
\diagram{\begin{tikzcd}[ampersand replacement=\&]
 \& B_{D_1^*}\arrow[dl]\arrow[dr] \& \& \\
 B_U \& \& B_D\arrow[r,leftrightarrow,"\unsure{\sim}"] \& B_S
\end{tikzcd}}
}

% Lignes 17--24.
Je m'aperçois qu'il me faut revenir sur les notations des \deleted{\illegible{1}} \inserted{divers} topos associés à une telle donnée. Mais je vais me guider sur la situation des n$^\circ$ .-- où on a un schéma régulier \inserted{$X$} de dim 1, un sous-schéma fermé \inserted{$S$} de dim 0, d'où $U=X-S$ --- dans ce cas \deleted{\illegible{2}} $\unsure{S_{\text{ét}}|X_{\text{ét}}}$ \textcircled{ne} peut se reconstituer à partir de $B_{D^*}$ comme $\unsure{B_{\pi_0 B_{D^*}}|B_I}$
\mspage{160}


\scanpage{162}
% Lignes 1--7.
avec $I=\pi_0(B_{D^*})$, il faut tenir compte des corps résiduels $k(s)$ $(s\in S)$ i.e. des groupes de Galois $\operatorname{Gal}(\overline{k(s)}/k(s))$. Donc il y a lieu de revenir à une situation de départ (qui est adaptée au cas \inserted{arithmético-}géométrique) de morphismes de topos \uuline{multigaloisiens}
\[
 (*)\qquad
 \diagram{\begin{tikzcd}[ampersand replacement=\&]
 \& B_{D^*}\arrow[dl,"\rho"']\arrow[dr,"\sigma"] \& \\
 B_U \& \& B_S
 \end{tikzcd}}
\]
% Lignes 8--14.
(Attention, on écrit $B_S$, non $B_D$, qui aura un autre sens), \uline{où} $\sigma$ \uline{induit un iso sur les} $\Pi_0$ [et est surjectif sur les $\operatorname{Hom}$]. Pour l'instant, on ne va faire aucune hypothèse particulière sur cette situation, qui pour $B_U$ connexe, et en termes des choix
% Lignes 15--23 : crochet et parenthèses tels que visibles.
[\deleted{d'objets $\widetilde D_i^*\in\unsure{D}$} de rev. universels $\widetilde D_i^*$ des comp. conn. $D_i^*$ de $D^*$, et d'un rev universel $\widetilde U$ de $B_U$, et de l'isom. entre les $\rho_!(\widetilde D_i^*)$ et $\widetilde U$) s'explicite par la donnée d'un groupe \inserted{$\mathcal E$ (ou} $\Pi$) $(=\operatorname{Aut}(\widetilde U))$ \deleted{\unsure{et}} \inserted{$\mathcal E_i$ (ou} $\Pi_i$ $(=\operatorname{Aut}(\widetilde D_i^*))$ \deleted{\illegible{2}} $(i\in I)$, $\Gamma_i$ $(i\in I)$
\mspage{161}


\scanpage{163}
\mspage{81}
% Lignes 1--3 et premier diagramme.
et des homomorphismes de groupes
\diagram{\begin{tikzcd}[ampersand replacement=\&,row sep=2em,column sep=2.5em]
\& \unsure{\mathcal E_i} \arrow[dl,no head,"\rho_i"'] \arrow[dr,"\sigma_i"] \& \\
\unsure{\mathcal E} \& \& \Gamma_i
\end{tikzcd}}
avec les $\sigma_i$ surjectifs. \deleted{\unsure{et}}
% Lignes 3--10.
[Quand il y a un corps de base $K$ pour la situation géométrique, alors ds la description \unsure{toposienne}, \inserted{\unsure{posant} $e=\operatorname{Spec}K$ et} désignant par $B_e$ le topos étale \inserted{$e_{\text{ét}}$} de $e$, i.e. $B_\Gamma$ \unsure{où|si} $\Gamma=\operatorname{Gal}(\overline K/K)$, le diagramme $*$ s'insère dans
\diagram{\begin{tikzcd}[ampersand replacement=\&]
 \& B_{D^*}\arrow[dl,"\rho"']\arrow[dr,"\sigma"] \& \\
 B_U\arrow[dr,"\varphi"'] \& \& B_S\arrow[dl,"\psi"] \\
 \& B_e \&
\end{tikzcd}},
avec donnée de commutativité pour le carré envisagé\unsure{\ldots\ comme envisagé} au début de ces notes].

% Lignes 11--18 et dernier diagramme.
En termes de $*$, on construit par ``recollement'' de $B_U$ et \inserted{de} $B_S$ (via le foncteur de recollement $\sigma_*\rho^*$) un topos mixte, qui n'est plus en général multigaloisien, noté précédemment $\unsure{D_{X,U}|B_{X,U}}$, et que \inserted{\unsure{je préfère}} maintenant \deleted{\illegible{1}} \unsure{noter|notés} $B_{X;S}$ \deleted{\illegible{2}} \inserted{[pour rappeler qu'il s'agit de faisceaux sur $X$, mais n'ayant de singularités que sur $S$]} \deleted{\illegible{1} s'insère dans un}. Il s'insère dans un diagramme de topos
\diagram{\begin{tikzcd}[ampersand replacement=\&,row sep=2em,column sep=2.5em]
\& B_{D^*} \arrow[dl] \arrow[dr] \& \\
B_U \arrow[dr,""{name=xl,description}] \& \& B_S \arrow[dl,""{name=xr,description}] \\
\& B_{X;S} \&
\arrow[from=xr,to=xl,bend right=25,"\alpha_X"]
\end{tikzcd}}
avec une flèche
\mspage{162}


\scanpage{164}
% Lignes 1--5.
``de commutation'' $\alpha$ qui n'est telle que par abus de langage --- ce n'est pas un iso \deleted{de \unsure{foncteurs}} mais un morphisme de foncteurs sans plus
\[
 \sigma^*\psi^*\deleted{\unsure{u}}\xrightarrow{\alpha_X}\rho^*\varphi^*.
\]
% Lignes 6--11.
De la $\overline{\mathrm m}$ façon, recollant $D^*$ et $S$ via $\sigma_*$ (i.e. remplaçant $B_U$ par $B_{D^*}$ dans la construction précédente de $\unsure{D_{X;S}|B_{X;S}}$), on trouve un topos, pas multigaloisien en général, noté $B_{D;S}$. Dans le \deleted{\unsure{présent}} modèle géométrique où on a $X,S$ comme dessus,
% Lignes 12--18.
$B_X$ correspond aux faisceaux sur $X$ qui sont \unsure{ess.|ens.|en.} loc. constants sur $U$ (et sur $S$, où ils n'ont pas de mérite) i.e. sur $U$ proviennent par l'image inverse \deleted{d'un faisceau} \inserted{par} $U\longrightarrow B_{\Pi_1 U}$ d'un faisceau sur $B_{\Pi_1 U}$.

% Lignes 19--25.
et de $\overline{\mathrm m}$ $B_{D;S}$ correspond aux faisceaux sur $D=\coprod_i\operatorname{Spec}(\mathcal O_i=\text{hensélisé de }\mathcal O_{X,\unsure{i|s}})$ \deleted{qui sont} qui sont loc constants sur $D^*=D-S$ (et sur $S$, sans mérite!) mais avec l'hyp. de dim faite -- --- \unsure{en|on} fait
\[
 B_{D;S}\simeq D_{\text{ét}}
\]
($D_i$ n'a que 2 pts, $D_i^*=\{\eta_i\}$ --- ). Ainsi $B_{D;S}$ \unsure{$(\simeq D_{\text{ét}})$}
\mspage{163}


\scanpage{165}
\mspage{82}
% Lignes 1--2 et diagramme supérieur.
(quand on part de $X,S$) s'insère dans un triangle de morphismes de topos
\diagram{\begin{tikzcd}[ampersand replacement=\&,row sep=2em,column sep=2.5em]
\& B_{D^*} \arrow[dl,equals,""{name=deq,description}] \arrow[dr,"\sigma",""{name=dsig,description}] \& \\
B_{D^*} \arrow[dr,"\varphi_D"'] \& \& B_S \arrow[dl,"\psi_D"] \\
\& B_{D;S} \&
\arrow[from=dsig,to=deq,bend left=22,"\alpha_D"']
\end{tikzcd}}
\margin{\uline{NB} dans le cas géométrique, \deleted{\illegible{1}} \unsure{trois} topos de ce diagramme sont des topos étales ($D^*_{\text{ét}},D_{\text{ét}},S_{\text{ét}}$) et les flèches $\varphi_D$ et $\psi_D$ correspondent à \unsure{des morph.} \unsure{mais non} \illegible{1} de topos)}
% Lignes 3--4.
$\alpha_D$, comme $\alpha_X$ ci-dessus, n'est qu'un vulgaire homom. (pas iso en général).

% Lignes 5--14.
\deleted{On trouve donc un \unsure{to.}}
Quand on parle du morphisme canonique de $B_{D^*}$ dans $B_{D;S}$ c'est \inserted{de} $\varphi_D$ (et non de $\psi_D\sigma$) qu'il s'agit --- dans le cas géométrique on a $B_{D^*}\simeq D^*_{\text{ét}}$, et $\varphi_D$ correspond à l'inclusion de schémas
\[
 D^*\unsure{\supset|=}D\setminus S\longrightarrow D
\]
\inserted{et $\psi_D:B_S\longrightarrow B_{D;S}$ à l'inclusion de schémas $S\longrightarrow D$}
alors que $\sigma$ ni $\psi_D\sigma$ ne correspondent en général à des morphismes de schémas.

% Lignes 15--16 et diagramme inférieur entièrement barré.
On obtient ainsi un diagramme de topos
\deleted{\diagram{\begin{tikzcd}[ampersand replacement=\&,row sep=5em,column sep=8em]
\& B_{D^*}
\arrow[dl,"\unsure{\rho_X}"',"\deleted{\unsure{B_{D^*}}}"{description,pos=0.56},"\deleted{\unsure{\varphi_D}}"{pos=0.25}]
\arrow[dl,bend right=38,"\deleted{\unsure{\illegible{1}}}"']
\arrow[dr,"\sigma"{xshift=0.65em,yshift=0.55em}]
\arrow[d] \& \\
B_U \arrow[dr,"\varphi"',""{name=oldl,description}] \& B_D \arrow[d] \& B_S \arrow[l] \arrow[dl,"\psi",""{name=oldr,description}] \\
\& B_X \&
\arrow[from=2-1,to=2-2,phantom,"\text{comm}"{description,yshift=1.3em}]
\arrow[from=2-2,to=2-3,phantom,"{\begin{gathered}\text{pseudo-}\\[-0.4ex]\text{comm}\end{gathered}}"{description,xshift=-1.5em,yshift=2.2em}]
\arrow[from=2-2,to=2-3,phantom,"\text{comm}"{description,xshift=-.8em,yshift=-1.3em}]
\arrow[from=oldr,to=oldl,bend right=20,"\deleted{\unsure{\illegible{1}}}"{pos=.8,xshift=-1.5em,yshift=-.2em}]
\end{tikzcd}}
où dans le cas géométrique}
\mspage{164}


\scanpage{166}
% Diagramme supérieur (**).
\[
 (**)\qquad
 \diagram{\begin{tikzcd}[ampersand replacement=\&, column sep=large,row sep=small]
 \& \& \uline{B}_S\arrow[dl]\arrow[dd,"{\Vert}" description] \\
 \uline{B}_{D^*}\arrow[r]\arrow[dd]\arrow[urr] \& B_{D;S}\arrow[dd] \& \\
 \& \& \uline{B}_S\arrow[dl] \\
 \uline{B}_{\unsure{U}}\arrow[r] \& B_{X;S} \&
 \end{tikzcd}}
\]
\margin{Attention, le triangle n'est pas ess. commutatif, mais \unsure{on a} pseudo-commutativité\ldots}
% Lignes 1--5.
où les trois topos soulignés $\uline{B}_U$, $\uline{B}_S$ et $\uline{B}_{D^*}$ sont multigaloisiens, et $B_{D;S}$ et $B_{X;S}$ sont composites (obtenus par recollement de deux topos multigaloisiens).
\inserted{Sauf erreur les deux carrés sont [\uline{non} seulement ess$^t$. comm. mais aussi] 2-cartésiens (dans la 2-cat. des topos).}

% Lignes 6--12 et diagramme médian.
Passant aux $B_{\Pi_1}$? des topos envisagés (correspondants aux objets loc. constants sur ces topos) on trouve \inserted{un topos} $B_{\deleted{\Pi_1}X}$ \inserted{$=\unsure{B_{\Pi_1 B_{X,S}}|B_{\Pi_1 B_{X;S}}}$} comme somme amalgamée de topos dans le \deleted{$B_{\Pi_1 X}$} diagramme
\diagram{\begin{tikzcd}[ampersand replacement=\&,row sep=2em,column sep=2.5em]
\& B_{D^*} \arrow[dl] \arrow[dr] \& \\
B_U \arrow[dr,""{name=pl,description}] \& \& B_S \arrow[dl,no head,""{name=pr,description}] \\
\& B_{\deleted{\unsure{\Pi_1}}X} \&
\arrow[from=pr,to=pl,bend right=25,"\sim"]
\end{tikzcd}}
% Lignes 13--19.
(ce carré est bel et bien essentiellement commutatif) et de $\overline{\mathrm m}$ (remplaçant $B_U$ par $B_{D^*}$) un $B_{\deleted{\Pi_1}D}$, qui est cependant (par $B_S\longrightarrow B_{\deleted{\Pi_1}D}$) isom. \inserted{(plutôt: \unsure{équiv.|égaux})} à $B_S$. Le diagramme $(**)$ devient alors
\mspage{165}


\scanpage{167}
\mspage{83}
% Diagramme supérieur (***).
\[
 (***)\qquad
 \diagram{\begin{tikzcd}[ampersand replacement=\&, column sep=large,row sep=small]
 \& \& B_S\;(=B_{\Pi_1 S})\arrow[dl,"\sim"']\arrow[dd,"{\Vert}" description] \\
 B_{D^*}\;(=B_{\Pi_1 D^*})\arrow[r]\arrow[dd]\arrow[urr,"\sigma"] \& B_{\deleted{\Pi_1}D}\arrow[dd] \& \\
 \& \& B_S\arrow[dl] \\
 B_U\;(=B_{\Pi_1 U})\arrow[r] \& B_{\deleted{\Pi_1}X} \&
 \end{tikzcd}}
\]
% Lignes 1--4.
où cette fois-ci le triangle sup. est bien ess. surjectif (c'est bien comme ça qu'on a défini $\sigma$, dans le cas géométrique!) et on a un morphisme de
% Lignes 5--12.
$(**)$ dans $(***)$, qui \deleted{\unsure{est}} pour les flèches qui \unsure{n'sont|sont} pas ess. des \inserted{\unsure{(semi-}}identités), $B_{D;S}\longrightarrow B_{\deleted{\Pi_1}D}$ et $B_{X;S}\longrightarrow B_{\deleted{\Pi_1}X}$, ont une nette tendance d'être ``acycliques'' i.e. d'induire des iso. sur la cohomologie (il faudrait voir \illegible{1} ce point). Enfin,
% Lignes 13--15 et diagramme inférieur (****).
dans le cas géométrique, on a \inserted{\unsure{des}} diagrammes \unsure{analogues|canoniques} de morphismes de topos étales \deleted{\illegible{1}}
\[
 (****)\qquad
 \diagram{\begin{tikzcd}[ampersand replacement=\&, column sep=large,row sep=small]
 \& \& S_{\text{ét}}\arrow[dl]\arrow[dd,equal] \\
 D^*_{\text{ét}}\arrow[r]\arrow[dd]\arrow[urr] \& D_{\text{ét}}\arrow[dd] \& \\
 \& \& S_{\text{ét}}\arrow[dl] \\
 U_{\text{ét}} \& X_{\text{ét}}\arrow[l] \&
 \end{tikzcd}}
\]
\margin{triangle sup. pas ess. comm.}
\mspage{166}


\scanpage{168}
% Lignes 1--2.
où toutes les flèches sauf $D^*_{\text{ét}}\longrightarrow S_{\text{ét}}$ sont induites par des morphismes de topos.
% Lignes 3--11.
Ici $(****)$ s'envoie dans le diagramme homologue $(**)$, en induisant des iso. de topos pour $D^*,D,S$, et pour $U_{\text{ét}}\longrightarrow B_U$ et $X_{\text{ét}}\longrightarrow X$ induisant des morphismes qui ont moins tendance d'être acycliques, mais qui le sont \deleted{\unsure{cependant}} quand $\overline{\mathrm m}$ dans des cas importants, rappelés dans une section antérieure\ldots

% Lignes 12--15 et diagramme.
Les constructions de $(**)$ et $(***)$ en termes du diagramme de départ
\diagram{\begin{tikzcd}[ampersand replacement=\&]
 \& B_{D^*}\arrow[dl]\arrow[dr] \& \\
 B_U \& \& B_S
\end{tikzcd}}
sont purement formelles, et indépendantes de \unsure{tte} hypothèse.
% Lignes 16--22.
La construction d'un $B_{\deleted{\Pi_1}X}$ multigaloisien, comme somme amalgamée, peut être interprété comme la traduction (au niveau des groupoïdes fondamentaux) d'une opération de ``bouchage de trous''.
\deleted{Au niveau} Dans le contexte calculatoire
\mspage{167}


\scanpage{169}
\mspage{84}
% Ligne 1 et diagramme de groupes.
(avec choix des $\widetilde U,\widetilde U_i$, $\rho_!(\widetilde U_i)\simeq\widetilde U$) on a
\diagram{\begin{tikzcd}[ampersand replacement=\&]
 \& \mathcal E_i\arrow[dl,"\rho_i"']\arrow[dr,"\sigma_i"] \& \\
 \pi \& \& \Gamma_i
\end{tikzcd}},
% Lignes 2--8.
posant $\widetilde X=\text{image de }\widetilde U$ par $i_!$ $(i:B_U\longrightarrow B_{\Pi_1 X})$ $B_{\deleted{\Pi_1}X}$ est décrit en termes de ce $\widetilde U$ comme le \deleted{topos des $\Pi$-ens} classifiant $B_{\pi_X}$, où $\pi_X=\pi_1(X;\widetilde X)$ \deleted{est défini} \inserted{se calcule} comme quotient de $\pi$ par le ss-groupe invariant engendré par les $\rho_i(L_i)$, où $L_i=\operatorname{Ker}\sigma_i\subset\mathcal E_i$.

% Lignes 9--16.
Si on se donne un sommande direct $B_{D_0^*}$ de $B_{D^*}$ (correspondant à une partie $I_0$ de $I=\pi_0(B_{D^*})$) on trouve de $\overline{\mathrm m}$ une somme amalgamée de $B_U$ et de $B_{S_0}$ sous $B_{D_0^*}$, noté \deleted{$\unsure{B_{\Pi_1(U\cup S_0)}}$} \inserted{$B_{U\cup S_1}\longleftarrow B_U$}, qui se visualise comme un bouchage partiel de trous, interprété au niveau des groupoïdes fondamentaux. Dans le cas géom, si on pose $I=I_0\amalg I_1$, i.e. $S=S_0\amalg S_1$, on
\mspage{168}


\scanpage{170}
% Lignes 1--2.
peut interpréter ce topos comme $B_{\Pi_1 U_1}$, où $U_1=X\setminus\unsure{S_1}$. \deleted{Bien sûr,}

% Lignes 3--6 et premier petit diagramme.
Bien sûr, on a un hom de diagrammes cocartésiens de topos relatif à
\diagram{\begin{tikzcd}[ampersand replacement=\&]
 \& B_{D_0^*}\arrow[dl]\arrow[dr] \& \\
 B_U \& \& B_{S_0}
\end{tikzcd}}
\deleted{\unsure{de dessous}} celui relatif à $B_{D^*}$ \deleted{\unsure{la situation}} s'envoyant dans $B_U$, et \deleted{$\unsure{B_{S_0}}$} $B_S$
% Grand diagramme central, avec surcharges et traits peu lisibles.
\diagram{\begin{tikzcd}[ampersand replacement=\&,row sep=3em,column sep=1.3em]
[B_{D_1^*}] \arrow[dr,dotted] \& \& B_{D_0^*} \arrow[dl] \arrow[dr] \arrow[rrr] \& \& \& B_{D^*} \arrow[dl,phantom,""{name=diagupper,description,pos=0.68,inner sep=0pt},""{name=diaglower,description,pos=0.86,inner sep=0pt}] \arrow[dr] \& \\
\& |[alias=leftU]| B_U \arrow[dr] \arrow[rrr,phantom,""{name=ugapleft,description,pos=0.39},""{name=ugapright,description,pos=0.76},""{name=utip,description,pos=0.88},""{name=uheadstart,description,pos=0.93,inner sep=0pt},""{name=uheadend,description,pos=0.97,inner sep=0pt},""{name=udashleft,description,pos=0.74,yshift=2pt,inner sep=0pt},""{name=udashright,description,pos=0.80,yshift=2pt,inner sep=0pt}] \& \& |[yshift=1.5em]| B_{S_0} \arrow[dl] \arrow[rrr] \& |[alias=rightU]| B_U \arrow[dr] \& \& |[yshift=1.5em]| B_S \arrow[dl] \\
\& \& B_{\deleted{\unsure{\Pi_1}}(\underbrace{U\cup S_0}_{U_1})} \arrow[rrr,no head] \& \& \&
\begin{gathered}
B_{\deleted{\unsure{\Pi_1}}(\underbrace{U\cup S}_{X})}\,\unsure{'}\\[-0.2ex]
\deleted{\unsure{\illegible{1}}\;B_{\unsure{\Pi_1}X}}
\end{gathered} \&
\arrow[from=1-6,to=diagupper,no head,shorten >=0pt]
\arrow[from=diaglower,to=rightU,no head,shorten <=0pt]
\arrow[from=leftU,to=ugapleft]
\arrow[from=ugapright,to=utip,no head]
\arrow[from=uheadstart,to=uheadend,shorten <=0pt,shorten >=0pt]
\arrow[from=udashleft,to=udashright,no head,shorten <=0pt,shorten >=0pt]
\end{tikzcd}}
,
% Lignes 7--9.
et d'autre part on a \inserted{un composé} $B_{D_1^*}\longrightarrow B_U\longrightarrow B_{\deleted{\Pi_1}U_1}$ qui avec $\unsure{B_{D_1^*}}\longrightarrow B_{S_1}$ \deleted{peut être utilisé} \inserted{donne un diagramme} de topos
\diagram{\begin{tikzcd}[ampersand replacement=\&]
 \& B_{D_1^*}\arrow[dl]\arrow[dr] \& \\
 B_{\deleted{\Pi_1}U_1} \& \& B_{S_1}
\end{tikzcd}}
% Lignes 10--15 (phrase continuée hors du lot).
de $\overline{\mathrm m}$ type qu'au début, qu'on peut utiliser pour construire encore une \deleted{\unsure{donc|dans}} somme amalgamée. Et il est évident que celle-ci est can. équivalente à $B_{\deleted{\Pi_1}X}$, la somme amalgamée corres-
\mspage{169}
